% amplifier-bundle-research: IMRAD LaTeX skeleton
% Target: arXiv preprint (default). Switch document class for venue-specific formatting.
% Usage: venue-formatter agent populates this skeleton during /publish mode.

\documentclass[11pt]{article}

% === Essential package stack ===
\usepackage[utf8]{inputenc}
\usepackage[T1]{fontenc}
\usepackage{amsmath,amssymb,mathtools}  % math
\usepackage{graphicx}                    % figures
\usepackage{booktabs}                    % tables (never use \hline with vertical rules)
\usepackage{microtype}                   % typography
\usepackage[numbers]{natbib}             % citations — switch to author-year for NeurIPS/ACL
\usepackage{hyperref}                    % clickable refs (load second-to-last)
\usepackage[capitalize,noabbrev]{cleveref} % smart refs (load last)

% === Optional but recommended ===
\usepackage{algorithm,algorithmic}       % algorithms
\usepackage{subcaption}                  % subfigures
\usepackage{xcolor}                      % colored text (for draft annotations)
\usepackage{enumitem}                    % list control
\usepackage{url}                         % URL formatting

% === Draft-mode annotations (remove before submission) ===
\newcommand{\todo}[1]{\textcolor{red}{\textbf{TODO:} #1}}
\newcommand{\exploratory}[1]{\textcolor{orange}{\textbf{[EXPLORATORY]} #1}}
\newcommand{\prereg}[1]{\textcolor{blue}{\textbf{[PRE-REG]} #1}}

% === Document metadata ===
\title{%TITLE%}
\author{%AUTHORS%}
\date{%DATE%}

\begin{document}

\maketitle

% === Abstract ===
% 5-component structure: background, gap, approach, results, implications
\begin{abstract}
%ABSTRACT_BACKGROUND%
%ABSTRACT_GAP%
%ABSTRACT_APPROACH%
%ABSTRACT_RESULTS%
%ABSTRACT_IMPLICATIONS%
\end{abstract}

% === Keywords (if required by venue) ===
% \textbf{Keywords:} %KEYWORDS%

% === 1. Introduction ===
\section{Introduction}
\label{sec:introduction}

% Paragraph 1: Broad context — why this area matters
%INTRO_CONTEXT%

% Paragraph 2: Specific problem — what gap exists
%INTRO_GAP%

% Paragraph 3: This work — what we do and why it matters
%INTRO_CONTRIBUTION%

% Paragraph 4: Results preview — what we found (1-2 sentences)
%INTRO_RESULTS_PREVIEW%

% Paragraph 5: Paper structure (optional — skip for short papers)
%INTRO_STRUCTURE%

% === 2. Related Work ===
\section{Related Work}
\label{sec:related}

%RELATED_WORK%

% === 3. Methods ===
\section{Methods}
\label{sec:methods}

% Pre-registration reference (if applicable)
% This study was pre-registered at [URL/hash]. The pre-registration
% specifies hypotheses, analysis plan, and stopping rules.

%METHODS_CONTENT%

\subsection{Experimental Setup}
\label{sec:setup}
%METHODS_SETUP%

\subsection{Evaluation Metrics}
\label{sec:metrics}
%METHODS_METRICS%

% === 4. Results ===
\section{Results}
\label{sec:results}

% Confirmatory results first, then exploratory (labeled)
%RESULTS_CONFIRMATORY%

%RESULTS_EXPLORATORY%

% === 5. Discussion ===
\section{Discussion}
\label{sec:discussion}

%DISCUSSION_CONTENT%

\subsection{Limitations}
\label{sec:limitations}
% Be specific — not "our study has limitations" but exactly what they are
%LIMITATIONS%

% === 6. Conclusion ===
\section{Conclusion}
\label{sec:conclusion}
%CONCLUSION%

% === Acknowledgments ===
% \section*{Acknowledgments}
% %ACKNOWLEDGMENTS%

% === Ethics Statement (if required) ===
% \section*{Ethics Statement}
% %ETHICS%

% === Reproducibility Statement ===
% \section*{Reproducibility Statement}
% Code: %CODE_URL%
% Data: %DATA_URL%
% Environment: %ENVIRONMENT%
% Seeds: %SEEDS%

% === Author Contributions (CRediT taxonomy) ===
% \section*{Author Contributions}
% %CREDIT_CONTRIBUTIONS%

% === Conflict of Interest ===
% \section*{Declaration of Competing Interests}
% %COI%

% === Data Availability ===
% \section*{Data Availability Statement}
% %DATA_AVAILABILITY%

% === Bibliography ===
\bibliographystyle{plainnat}  % switch per venue: abbrvnat, ieeetr, acl_natbib
\bibliography{references}

% === Appendix (optional) ===
% \appendix
% \section{Supplementary Material}
% \label{app:supplementary}
% %APPENDIX%

\end{document}
